%=====================================================================
\section{Conclusion}
\label{sec:conclusion}
%=====================================================================

Embedding-level erasure can only remove the concept features it finds.
We show that \ember{} often misses them because concepts are quiet directions of the embedding space, overshadowed by generic topics, and because fixed feature widths and summary-based judging lose what remains.
We address this with \cinder{}, which sharpens the embedding matrix toward the concept, lets each feature find its width, verifies features both bottom-up and top-down, and erases them relative to ordinary text.
Sharpening against a multilingual background also turns embedding-level erasure, so far limited to English, into a single fit for eight languages that matches erasing each language separately with an eighth of the judge calls.
We also find that what the model's first layer reads in an edit is dominated by one direction shared by all tokens, which separates what an edit removes from a concept from what it removes from every word.
Evaluating \cinder{}, we further show that the standard protocol rewards destroying the edited tokens, and propose a collateral-aware evaluation that charges for it.
Under this evaluation, embedding erasure appears as what it is, a way to block access to a concept through its words, and within that trade-off \cinder{} reaches more concepts and, at moderate edit size, removes more.
These findings suggest that improving embedding erasure starts with finding the concept, and that its evaluation must look at the words it edits.
